% Portada de C9: una petición a la API, la respuesta con su uso de tokens y tres cuidados.
\documentclass[tikz,border=0pt]{standalone}
\usepackage{resonancia}
\begin{document}
\begin{tikzpicture}
  \clip[rounded corners=6pt] (0,0) rectangle (12,3.6);
  \shade[top color=rGris!22, bottom color=white] (0,0) rectangle (12,3.6);
  \node[caja, draw=rAzul, fill=white, minimum width=2.6cm, minimum height=1.0cm, font=\sffamily\small, text=rAzul] (prog) at (1.8,2.5) {tu programa};
  \node[caja, draw=rVioleta, fill=white, minimum width=2.6cm, minimum height=1.0cm, font=\sffamily\small, text=rVioleta] (api) at (10.2,2.5) {API de Claude};
  \draw[trazo, color=rAzul, -{Stealth}] (prog.north east) .. controls (5,3.6) and (7,3.6) .. (api.north west);
  \node[etiqueta, color=rAzul] at (6,3.35) {mensajes, max\_tokens};
  \draw[trazo, color=rVioleta, -{Stealth}] (api.south west) .. controls (7,1.5) and (5,1.5) .. (prog.south east);
  \node[etiqueta, color=rVioleta] at (6,1.65) {respuesta, stop\_reason, usage};
  \foreach \x/\t/\c in {2.2/clave en el entorno/rNaranja, 6/coste desde usage/rVerde, 9.8/errores con criterio/rGris} {
    \node[caja, draw=\c, fill=white, minimum width=3.1cm, minimum height=0.6cm, inner sep=2pt, font=\sffamily\scriptsize, text=\c] at (\x,0.5) {\t};
  }
\end{tikzpicture}
\end{document}
